\documentclass[11pt,a4paper,openany]{book}

\usepackage[T1]{fontenc}
\usepackage[utf8]{inputenc}
\usepackage[italian]{babel}
\usepackage{lmodern}
\usepackage{microtype}
\usepackage[a4paper,margin=2.4cm,headheight=15pt]{geometry}
\usepackage{needspace}
\usepackage{xcolor}
\usepackage{amsmath,amssymb}
\usepackage{booktabs}
\usepackage{enumitem}
\usepackage{listings}
\usepackage[most]{tcolorbox}
\usepackage{fancyhdr}
\usepackage{titlesec}
\usepackage{hyperref}

\definecolor{Ink}{HTML}{172033}
\definecolor{Slate}{HTML}{526075}
\definecolor{Paper}{HTML}{F7F8FA}
\definecolor{Accent}{HTML}{2457C5}
\definecolor{Bad}{HTML}{A83232}
\definecolor{BadBack}{HTML}{FFF2F1}
\definecolor{Good}{HTML}{18724B}
\definecolor{GoodBack}{HTML}{EFFAF4}
\definecolor{Note}{HTML}{9A6700}
\definecolor{NoteBack}{HTML}{FFF8DF}
\definecolor{CodeBack}{HTML}{F3F5F8}

\hypersetup{
  colorlinks=true,
  linkcolor=Accent,
  urlcolor=Accent,
  pdftitle={Dispensa progressiva di algoritmi e complessità},
  pdfauthor={Giuseppe Sinatra}
}

\setlength{\parindent}{0pt}
\setlength{\parskip}{0.65em}
\setlist[itemize]{topsep=0.25em,itemsep=0.2em}
\setlist[enumerate]{topsep=0.25em,itemsep=0.2em}

\titleformat{\chapter}[display]
  {\normalfont\bfseries\color{Ink}}
  {\filleft\Large\color{Accent}\chaptername\ \thechapter}
  {0.6ex}
  {\titlerule[1.1pt]\vspace{1.2ex}\Huge}
\titleformat{\section}
  {\Large\bfseries\color{Ink}}{\thesection}{0.7em}{}
\titleformat{\subsection}
  {\large\bfseries\color{Ink}}{\thesubsection}{0.7em}{}

\pagestyle{fancy}
\fancyhf{}
\fancyhead[L]{\small\color{Slate}Coding Skill Tree}
\fancyhead[R]{\small\color{Slate}Dispensa progressiva}
\fancyfoot[C]{\small\color{Slate}\thepage}
\renewcommand{\headrulewidth}{0.3pt}
\let\cleardoublepage\clearpage

\lstdefinestyle{pythoncode}{
  language=Python,
  basicstyle=\ttfamily\small,
  keywordstyle=\bfseries\color{Accent},
  commentstyle=\itshape\color{Slate},
  stringstyle=\color{Good},
  showstringspaces=false,
  columns=fullflexible,
  keepspaces=true,
  breaklines=true,
  breakatwhitespace=false,
  tabsize=4,
  numbers=left,
  numberstyle=\scriptsize\color{Slate},
  numbersep=8pt,
  xleftmargin=1.2em
}

\newtcblisting{inefficientcode}[1]{
  enhanced,
  breakable,
  listing only,
  listing options={style=pythoncode},
  title={#1},
  colback=BadBack,
  colframe=Bad,
  coltitle=white,
  fonttitle=\bfseries,
  boxrule=0.9pt,
  arc=2pt,
  left=1mm,right=1mm,top=1mm,bottom=1mm
}

\newtcblisting{efficientcode}[1]{
  enhanced,
  breakable,
  listing only,
  listing options={style=pythoncode},
  title={#1},
  colback=GoodBack,
  colframe=Good,
  coltitle=white,
  fonttitle=\bfseries,
  boxrule=0.9pt,
  arc=2pt,
  left=1mm,right=1mm,top=1mm,bottom=1mm
}

\newtcolorbox{principlebox}[1]{
  enhanced,
  breakable,
  title={#1},
  colback=NoteBack,
  colframe=Note,
  coltitle=white,
  fonttitle=\bfseries,
  boxrule=0.9pt,
  arc=2pt
}

\newtcolorbox{definitionbox}[1]{
  enhanced,
  breakable,
  title={#1},
  colback=Paper,
  colframe=Accent,
  coltitle=white,
  fonttitle=\bfseries,
  boxrule=0.9pt,
  arc=2pt
}

\begin{document}

\frontmatter
\begin{titlepage}
  \centering
  \vspace*{2.2cm}
  {\Large\color{Accent}\bfseries CODING SKILL TREE\par}
  \vspace{1.5cm}
  {\Huge\color{Ink}\bfseries Dispensa progressiva di algoritmi e complessità\par}
  \vspace{0.7cm}
  {\Large\color{Slate}Dalle implementazioni non conformi agli invarianti efficienti\par}
  \vfill
  \begin{tcolorbox}[
    width=0.82\textwidth,
    colback=Paper,
    colframe=Accent,
    boxrule=0.8pt,
    arc=2pt
  ]
  \centering
  Manuale cumulativo per il ripasso offline dei pattern algoritmici emersi
  nelle quest settimanali e nelle relative review tecniche.
  \end{tcolorbox}
  \vfill
  {\large Giuseppe Sinatra\par}
  \vspace{0.3cm}
  {\color{Slate}Edizione progressiva - aggiornata al 5 agosto 2026\par}
\end{titlepage}

\chapter*{Nota metodologica}
\addcontentsline{toc}{chapter}{Nota metodologica}

Questa dispensa non è una raccolta generale di algoritmi. Ogni capitolo nasce
da un'evidenza osservabile: un pattern nuovo affrontato in assessment, una
soluzione non corretta oppure un'implementazione corretta sugli esempi ma non
conforme ai vincoli di scala.

L'analisi segue una struttura costante: modello del problema, approccio
iniziale, diagnosi, invariante dell'algoritmo efficiente, argomento di
correttezza, complessità ed elementi utili al riconoscimento del pattern.
I riquadri rossi conservano il valore diagnostico dell'implementazione
iniziale; i riquadri verdi presentano la versione di riferimento validata in
review.

\begin{principlebox}{Regola di lettura}
Il codice rosso non è necessariamente errato negli output. Può essere
semanticamente corretto e, al tempo stesso, inaccettabile rispetto ai limiti
computazionali del problema. La distinzione tra correttezza funzionale ed
efficienza asintotica è parte integrante della valutazione algoritmica.
\end{principlebox}

\tableofcontents

\mainmatter
\chapter{Predicati esatti e sliding window}

\begin{flushleft}
\textbf{Origine:} Week 001, esercizi 2 e 4\\
\textbf{Contesto:} aggregazione di record e picco di richieste temporali\\
\textbf{Esito della review:} un errore logico e un vincolo di scala non rispettato\\
\textbf{Pattern:} predicati esatti; sliding window a due indici
\end{flushleft}

\section{Predicati esatti nell'aggregazione di record}

\subsection{Modello, ipotesi e notazione}

Sia $R=(r_0,\ldots,r_{n-1})$ una sequenza di record. Ogni record può
contenere i campi \texttt{model}, \texttt{tokens} e \texttt{success}. Per
ogni nome di modello valido occorre calcolare numero di run, numero di
successi e somma dei token. La condizione sui successi è intenzionalmente
stretta:
\[
\operatorname{successo}(r_i) \iff r_i[\texttt{success}]\ \texttt{is True}.
\]
Il contratto non afferma che il campo contenga esclusivamente booleani;
pertanto non è lecito sostituire questa condizione con la truthiness.

\subsection{Approccio iniziale ed errore logico}

La submission costruisce correttamente un dizionario indicizzato per modello,
gestisce token mancanti o nulli e non modifica gli input. La riga che calcola
\texttt{successful\_run}, tuttavia, accetta ogni oggetto truthy. Ad esempio,
\texttt{1} e \texttt{"yes"} aumentano il contatore pur non essendo l'oggetto
booleano \texttt{True}. Si tratta di un errore di correttezza, non di una
semplice inefficienza.

Il codice seguente conserva integralmente la submission; sono omessi soltanto
gli spazi finali invisibili, privi di effetto sull'esecuzione, per esigenze di
presentazione tipografica.

\begin{inefficientcode}{Submission originale - errore logico sul predicato}
def summarize_model_runs(runs):
    """Aggrega run, successi e token per modello."""
    lookup = {}
    for record_or in runs:
        record=record_or.copy()
        model=record.get('model')
        tokens= 0 if record.get('tokens') is None else record.get('tokens')
        successful_run=1 if record.get("success") else 0
        if model is None:
            continue
        if model not in lookup:
            lookup[model]={"total_runs": 1, "successful_runs": successful_run, "total_tokens": tokens}
        else:
            lookup[model]["total_runs"]+=1
            lookup[model]["successful_runs"]+=successful_run
            lookup[model]["total_tokens"]+=tokens

    return lookup
\end{inefficientcode}

\subsection{Algoritmo corretto e invariante}

La struttura di aggregazione resta appropriata. È sufficiente inizializzare
le statistiche del modello, aggiornare esattamente una volta il numero di run
e tradurre letteralmente il predicato con \texttt{is True}.

\begin{principlebox}{Invariante dell'aggregazione}
Dopo avere elaborato i primi $i$ record, per ogni modello $m$ presente nel
dizionario le tre statistiche sono esattamente quelle dei record validi tra
$r_0$ e $r_{i-1}$ con modello $m$; in particolare,
\texttt{successful\_runs} conta soltanto i record il cui valore di
\texttt{success} è l'oggetto \texttt{True}.
\end{principlebox}

\begin{efficientcode}{Implementazione di riferimento - predicato conforme}
def summarize_model_runs(runs):
    """Aggrega run, successi e token per modello."""
    summary = {}
    for record in runs:
        model = record.get("model")
        if model is None:
            continue

        if model not in summary:
            summary[model] = {
                "total_runs": 0,
                "successful_runs": 0,
                "total_tokens": 0,
            }

        stats = summary[model]
        stats["total_runs"] += 1
        stats["successful_runs"] += int(record.get("success") is True)
        tokens = record.get("tokens")
        stats["total_tokens"] += 0 if tokens is None else tokens

    return summary
\end{efficientcode}

\subsection{Correttezza, complessità ed edge case}

L'invariante vale inizialmente perché il dizionario è vuoto. Ogni record senza
modello valido viene ignorato come richiesto. Per un record valido,
l'algoritmo aggiorna una sola voce: incrementa il numero di run, aggiunge un
successo se e solo se il predicato esatto è vero e somma zero oppure il numero
di token. L'invariante è quindi preservato; dopo $n$ iterazioni descrive
l'intero input e implica la correttezza del risultato.

Il tempo è $O(n)$ in media, assumendo accessi hash $O(1)$, e lo spazio
ausiliario è $O(k)$ per $k$ modelli distinti. I casi da verificare includono
lista vuota, modello mancante o \texttt{None}, nome di modello vuoto, token
mancanti, nulli o uguali a zero, e valori \texttt{success} quali
\texttt{False}, \texttt{None}, \texttt{0}, \texttt{1} e stringhe non vuote.

\begin{principlebox}{Segnale di riconoscimento}
Espressioni come ``soltanto quando'', ``esattamente'', ``è \texttt{True}'' o
``è \texttt{None}'' richiedono di distinguere identità, uguaglianza e
truthiness. Prima di codificare, trascrivere il predicato della specifica in
una condizione booleana senza allargarne il dominio.
\end{principlebox}

\section{Sliding window: modello del problema}

Sia
\[
T = (t_0,t_1,\ldots,t_{n-1})
\]
una sequenza di timestamp ordinata in modo non decrescente e sia $w>0$ la
durata della finestra. Si vuole determinare il massimo numero di eventi
contenuti in un intervallo semiaperto di ampiezza $w$:
\[
M = \max_{s}\left|\{t_i \in T \mid s \leq t_i < s+w\}\right|.
\]

La scelta dell'intervallo semiaperto è sostanziale: un evento posto
esattamente in $s+w$ non appartiene alla finestra che inizia in $s$.

\begin{definitionbox}{Precondizioni operative}
La soluzione lineare presuppone timestamp già ordinati. Se l'ordinamento non
è garantito dal contratto, occorre prima ordinare la sequenza, portando il
costo complessivo a $O(n\log n)$. La durata deve essere strettamente positiva;
in caso contrario la funzione solleva \texttt{ValueError}.
\end{definitionbox}

\Needspace{14cm}
\section{Sliding window: approccio iniziale}

L'implementazione presentata durante la quest seleziona ogni timestamp come
possibile inizio e, per ciascuno, visita nuovamente l'intera sequenza. Questa
strategia è concettualmente semplice e restituisce gli output attesi, ma
esegue fino a $n^2$ verifiche di appartenenza.

\begin{inefficientcode}{Implementazione iniziale - non conforme al vincolo di scala}
def max_requests_in_window(timestamps, window_seconds):
    """Restituisce il picco di richieste in una finestra temporale."""
    count_max=0
    if window_seconds <= 0:
        raise ValueError("La finestra dev'essere un valore positivo")
    for start in timestamps:
        count=0
        for input in timestamps:
            if start <= input < start+window_seconds:
                count+=1
        count_max=max(count, count_max)

    return count_max
\end{inefficientcode}

\subsection{Diagnosi tecnica}

Il problema non risiede nella condizione di appartenenza, che rappresenta
correttamente l'intervallo semiaperto. L'inefficienza nasce dalla perdita di
informazione tra due finestre consecutive: quando l'estremo destro avanza,
gli eventi ancora validi vengono contati nuovamente da zero.

Inoltre, il nome locale \texttt{input} oscura la funzione built-in omonima di
Python. Non altera il risultato in questo caso, ma riduce la chiarezza del
codice e può generare errori in estensioni successive.

Con un limite di $n=200\,000$, una crescita quadratica è incompatibile con un
coding assessment ordinario:
\[
n^2 = 40\,000\,000\,000
\]
controlli potenziali. Il superamento dei test visibili certifica soltanto
alcuni risultati, non il rispetto dei vincoli computazionali.

\clearpage
\thispagestyle{fancy}
\noindent\rule{0pt}{1cm}\par
\section{Sliding window: algoritmo efficiente}

Poiché i timestamp sono ordinati, gli eventi validi per un dato estremo
destro formano sempre un segmento contiguo della sequenza. È quindi
sufficiente mantenere due indici:

\begin{itemize}
  \item \texttt{right} introduce il nuovo evento corrente;
  \item \texttt{left} avanza soltanto quando l'evento più vecchio non rientra
        più nella finestra.
\end{itemize}

\begin{principlebox}{Invariante della finestra}
Dopo il ciclo \texttt{while}, tutti e soli gli eventi con indice compreso tra
\texttt{left} e \texttt{right} appartengono alla finestra semiaperta che
termina sul timestamp corrente. In particolare,
\[
t_{\texttt{right}}-t_{\texttt{left}} < w.
\]
L'indice \texttt{left} è il più piccolo indice che rende vera questa
disuguaglianza.
\end{principlebox}

\begin{efficientcode}{Implementazione di riferimento - complessità lineare}
def max_requests_in_window(timestamps, window_seconds):
    """Restituisce il picco di richieste in una finestra temporale."""
    if window_seconds <= 0:
        raise ValueError("La finestra dev'essere un valore positivo")

    left = 0
    best = 0

    for right, current_time in enumerate(timestamps):
        while current_time >= timestamps[left] + window_seconds:
            left += 1

        current_count = right - left + 1
        best = max(best, current_count)

    return best
\end{efficientcode}

\section{Sliding window: argomento di correttezza}

Si consideri un indice destro fissato. Prima dell'esecuzione del
\texttt{while}, la finestra può contenere timestamp troppo vecchi. Ogni
iterazione incrementa \texttt{left} finché il timestamp corrente non dista
meno di $w$ dal timestamp più vecchio conservato. Per ordinamento, tutti gli
elementi successivi a \texttt{left} soddisfano allora la stessa condizione.

Alla terminazione del ciclo, \texttt{left} è minimo: se esistesse un indice
valido più piccolo, il \texttt{while} si sarebbe fermato prima. Di
conseguenza, \texttt{right-left+1} è il massimo numero di eventi in una
finestra valida il cui evento più recente è $t_{\texttt{right}}$.
Considerando in successione ogni possibile estremo destro e conservando il
massimo, l'algoritmo restituisce $M$.

\clearpage
\thispagestyle{fancy}
\noindent\rule{0pt}{2cm}\par
\section{Sliding window: analisi di complessità}

\begin{center}
\begin{tabular}{@{}lll@{}}
\toprule
Strategia & Tempo & Spazio ausiliario \\
\midrule
Scansione annidata & $O(n^2)$ & $O(1)$ \\
Sliding window & $O(n)$ & $O(1)$ \\
Ordinamento + sliding window & $O(n\log n)$ & dipende dall'ordinamento \\
\bottomrule
\end{tabular}
\end{center}

Il ciclo \texttt{while} annidato non rende quadratica la seconda soluzione.
L'indice \texttt{left} non arretra mai: nel corso dell'intera esecuzione può
avanzare al massimo $n$ volte. Anche \texttt{right} percorre la sequenza una
sola volta; il lavoro totale è quindi lineare.

\section{Sliding window: edge case rilevanti}

\begin{itemize}
  \item \textbf{Sequenza vuota:} il ciclo non viene eseguito e il risultato è
        zero.
  \item \textbf{Timestamp duplicati:} restano simultaneamente nella finestra
        e vengono conteggiati come eventi distinti.
  \item \textbf{Evento sul bordo destro:} se la distanza è esattamente $w$,
        l'evento più vecchio viene escluso, coerentemente con $[s,s+w)$.
  \item \textbf{Durata non positiva:} il problema non definisce una finestra
        valida; la funzione segnala esplicitamente l'errore.
  \item \textbf{Input non ordinato:} l'invariante non è garantito; occorre
        ordinare oppure rifiutare l'input secondo il contratto.
\end{itemize}

\section{Sliding window: come riconoscere il pattern}

La sliding window è una candidata naturale quando il problema presenta
contemporaneamente:

\begin{itemize}
  \item una sequenza ordinata o ordinabile;
  \item un intervallo contiguo definito da una soglia;
  \item una statistica da massimizzare o minimizzare su tutti gli intervalli;
  \item finestre consecutive che condividono gran parte degli elementi.
\end{itemize}

\begin{principlebox}{Regola operativa per l'assessment}
Quando due finestre consecutive differiscono per pochi elementi, chiedersi
quale stato possa essere aggiornato incrementalmente. Se ogni elemento entra
ed esce dalla struttura al massimo una volta, è spesso possibile ottenere
una soluzione $O(n)$.
\end{principlebox}

\clearpage
\thispagestyle{fancy}
\noindent\rule{0pt}{2cm}\par
\section{Checklist di ripasso}

Prima di considerare concluso un problema analogo, verificare:

\begin{enumerate}
  \item se l'ordinamento è garantito oppure deve essere costruito;
  \item se gli estremi dell'intervallo sono inclusivi o esclusivi;
  \item quale invariante deve valere dopo l'avanzamento di \texttt{left};
  \item quante volte ciascun indice può avanzare complessivamente;
  \item se la complessità ottenuta è compatibile con la dimensione massima
        dichiarata.
\end{enumerate}

\begin{definitionbox}{Stato di apprendimento}
Il pattern è stato introdotto e analizzato, ma non è ancora masterizzato. La
mastery richiederà implementazioni lineari autonome, corrette e completate nel
timebox su più problemi con formulazioni differenti.
\end{definitionbox}


\backmatter
\chapter*{Indice dei pattern}
\addcontentsline{toc}{chapter}{Indice dei pattern}

\begin{tabular}{@{}lll@{}}
\toprule
Pattern & Prima introduzione & Stato al capitolo \\
\midrule
Sliding window a due indici & Week 001 & Introdotto, da consolidare \\
\bottomrule
\end{tabular}

\end{document}
